\section{财产犯罪 - 诈骗罪、职务侵占罪、普通罪名}
\begin{enumerate}
  \item 诈骗。财产犯罪里最难，最重要的。
        \begin{enumerate}
          \item \zd{非法占有目的 \yjt 欺骗行为 \yjt 产生 \hz 维持认知错误 \yjt 处分 \yjt 取得财物 \yjt 遭受财产损失}。
          \item 欺骗行为。
                \begin{enumerate}
                  \item 欺骗内容。事实判断可以，价值评价难以形成欺骗。
                        \begin{enumerate}
                          \item \al{交 1W 中介费，婚介所介绍美若天仙的女子给你}。
                        \end{enumerate}
                  \item 方式。不作为欺骗，隐瞒真相，如保险要告知自己真实情况。
                  \item 程度。足以使一般人相信。\al{我是秦始皇，打钱}。
                \end{enumerate}
          \item 人的要求。有意思自治能力。\al{骗小孩的钱、自动售卖机}。
          \item 基于认识错误处分财物。
                \begin{enumerate}
                  \item 客观处分行为。处分占有。\al{甲对柜台售货员说想试带手镯，带上跑掉}。
                  \item 主观处分意识。\zd{意识到财物的现实存在}。
                        \begin{enumerate}
                          \item \al{掉包案。把相机放到一箱饼干里去结账}。
                          \item \al{调换价签；缺斤少两}。
                        \end{enumerate}
                \end{enumerate}
          \item 损失。\zd{目的失败理论}。花冤枉钱。\al{甲买 200 元的柏浪涛刑法书，卖家寄送 200 元民法书}。
          \item 三角诈骗。看第三人\zd{是否有权处分的权利地位}。
                \begin{enumerate}
                  \item \al{甲来到乙家，对保姆谎称是干洗店老板，来取衣服}。
                  \item \al{甲骗店员乙去丙家院子取衣服}。
                  \item \al{地铁上乙问是不是甲的包，甲说是的，然后递给甲，甲随机下车，其实是丙的}。不构成。
                \end{enumerate}
        \end{enumerate}
  \item 职务侵占罪。\zd{A + B}。A，盗窃、侵占、诈骗。
        \begin{enumerate}
          \item B1 行为主体，单位职员，职务有一定要求，管理性。
          \item B2 行为对象，单位财物。
          \item B3（实质）利用职务便利。\al{财务总监去仓库偷东西}，形式利用。
        \end{enumerate}
  \item 故意毁坏财物罪。
        \begin{enumerate}
          \item 物理性毁损。
          \item （主客观）效用或功能丧失减损。\al{高速路撒钉子；乙往甲鱼塘投放过量蟹苗；往室友饭盒撒尿}。
          \item 主观效用标准。一般人价值观判断。
        \end{enumerate}
  \item 普通罪名。
        \begin{enumerate}
          \item 挪用资金罪。
                \begin{enumerate}
                  \item 超过 3 个月；未超 3 月 + 数额大 \hz 营利活动 \hz 非法活动。
                  \item 故意不归还 \yjt 职务侵占罪。
                  \item 进行非法活动构成其他罪的，数罪并罚。
                \end{enumerate}
          \item 挪用特定款物罪。
                \begin{enumerate}
                  \item 单位犯罪，但只处罚直接责任人。
                  \item 改变款物的专用用途。
                \end{enumerate}
          \item 拒不支付劳动报酬罪。
                \begin{enumerate}
                  \item 不作为犯。
                  \item 前提：经政府有关部门责令支付仍不支付。
                \end{enumerate}
        \end{enumerate}
\end{enumerate}
